\makeatletter
\def\input@path{{../styles/}{../../styles/}{../../../styles/}{../}{../../}{../../../}}
\makeatother

\documentclass{ee102_pset}
% macros.tex - Course meta information
\renewcommand{\course}{EE 102} % with a space
\renewcommand{\coursetitle}{Signal Processing and Linear Systems}
\renewcommand{\instructor}{Ayush Pandey}
\renewcommand{\student}{Name: }

\renewcommand{\semester}{Fall 2025}
\date{\semester} % this sets the LaTeX date field safely

% Problem set number
\renewcommand{\psetnum}{1}

% Release / due — use renewcommand because package provides empties
\renewcommand{\releasedate}{August 20, 2025}
\renewcommand{\duedate}{August 27, 2025}

\input{preamble.tex}

\author{\rule{3cm}{0.4pt}}
\submitdate{\rule{3cm}{0.4pt}}
\problemset{\shortstack{Problem Set \#8: Practice and Study Guide\\for Midterm Exam \#2}}
\renewcommand{\psetnum}{8}
\renewcommand{\releasedate}{October 19, 2026}
\renewcommand{\duedate}{Oct 28, 2026 at 11:59 PM}
\shorttitle{Problem Set \#8}
\hypersetup{pdftitle={EE 102 Problem Set 8: Practice and Study Guide for Midterm Exam 2}}

\begin{document}
This problem set is intentionally longer than usual to provide practice for
Midterm Exam \#2 on Wednesday, October 28, during class (1:40--2:55 PM).
Work through it before the exam. The submission deadline is
\textbf{Wednesday, October 28, at 11:59 PM}, after the exam, following our
usual submission schedule. The exam will be shorter than this practice set.

There are \textbf{100 points}. Show your derivations, label sketches, and
state any assumptions. Use angular frequency $\omega$ in radians per second.
Throughout this set, the Fourier transform convention is
\[
X(\omega)=\int_{-\infty}^{\infty}x(t)e^{-j\omega t}\,dt,
\qquad
x(t)=\frac{1}{2\pi}\int_{-\infty}^{\infty}X(\omega)e^{j\omega t}\,d\omega.
\]

\problem{1} \textbf{Complex exponentials and signal representations [10 points]}

\problempart[3 points] Write
\[
x(t)=6e^{-2t}\cos\left(4t+\frac{\pi}{3}\right)u(t)
\]
as $x(t)=\operatorname{Re}\{C e^{st}\}u(t)$. Find $C$ in rectangular and
polar form and find $s$. Explain which part of $s$ controls decay and which
part controls oscillation.

\problempart[3 points] Consider $e^{j3t}$, $\cos(3t)$, and $e^{-2t}u(t)$.
For each, determine whether it is periodic and whether its total energy is
finite. Which signals have a Fourier \emph{series} representation? Which
have a Fourier transform given by an ordinary convergent integral? Explain
how frequency-domain impulses represent the periodic examples.

\problempart[4 points] Use Euler's formula to express
$A\cos(\omega_0t+\phi)$ as two complex exponentials, where
$A>0$ and $\omega_0>0$. Give the coefficient of each exponential. Explain
why this real signal has contributions at both $+\omega_0$ and
$-\omega_0$, while $e^{j\omega_0t}$ has only one frequency component.

\vspace*{\fill}
\begin{center}[use additional pages if needed]\end{center}

\problem{2} \textbf{Circuit impulse responses and convolution [15 points]}
% Adapted from Fall 2025 HW8 Problem 1; input/output and initial conditions
% are now explicit. Circuit analysis is prerequisite knowledge.
Consider an ideal series RC circuit driven by an input voltage
$v_{\mathrm{in}}(t)$, with $R>0$, $C>0$, and output measured
\textbf{across the capacitor}. Let $\tau_C=RC$.
Assume zero capacitor voltage before the input is applied and use the
corresponding zero-state LTI model.

\problempart[4 points] Write down the RC impulse response $h_C(t)$ from
our earlier notes and briefly explain what an impulse response means.
State why the factor multiplying the exponential is needed; check that
the units of $h_C(t)$ make $(v_{\mathrm{in}}*h_C)(t)$ a voltage.

\problempart[6 points] For $v_{\mathrm{in}}(t)=V_0u(t)$, use convolution
with your impulse response to find the output voltage. Sketch it and
compare its value just after switching and as $t\to\infty$.

\problempart[5 points] Now let
$v_{\mathrm{in}}(t)=V_0[u(t)-u(t-D)]$, with $D>0$. Use your step-response
result and LTI properties to find the output voltage for all $t$.
Describe what happens at the two switching times.

\vspace*{\fill}
\begin{center}[use additional pages if needed]\end{center}

\problem{3} \textbf{Complex exponentials through LTI systems [10 points]}

\problempart[3 points] Starting from the convolution integral, show that
an LTI system with impulse response $h(t)$ responds to
$x(t)=e^{j\omega t}$, present for all time, according to
\[
y(t)=H(j\omega)e^{j\omega t},\qquad
H(j\omega)=\int_{-\infty}^{\infty}h(\tau)e^{-j\omega\tau}\,d\tau.
\]
Assume the integral exists, using sifting for any ideal impulse term.
Why is $e^{j\omega t}$ called an eigenfunction of an LTI system?

\problempart[3 points] Evaluate $H_C(j\omega)$ for
the impulse response in Problem 2. Write the
output to $A\cos(\omega t)$ in magnitude/phase form, assuming the input
has been present for all time and the circuit is in sinusoidal steady state.
Take $\omega>0$.

\problempart[4 points] For the RC circuit, apply the convolution integral
to the \emph{switched-on} input $x(t)=e^{j\omega t}u(t)$ and obtain its
zero-state output. Identify the transient term and show how your result
approaches the sinusoidal steady-state response. Explain why simply writing
$H_C(j\omega)e^{j\omega t}u(t)$ misses part of the response.

\vspace*{\fill}
\begin{center}[use additional pages if needed]\end{center}

\problem{4} \textbf{Fourier series and periodic filtering [10 points]}
% The pulse train replaces the stale cross-reference to last year's HW6.
Let $A>0$, $T>0$, and $\omega_0=2\pi/T$. Define the periodic pulse train
$q(t+T)=q(t)$ by one period:
\[
q(t)=\begin{cases}
A, & |t|<T/4,\\
0, & T/4<|t|<T/2.
\end{cases}
\]
Values at the isolated switching instants do not affect the coefficients.
Let $a_k$ denote its complex Fourier series coefficients.

\problempart[4 points] Sketch two periods and compute $a_0$ and $a_k$
for $k\ne0$ using the analysis equation. Check your result using the
symmetry of $q(t)$. Which harmonics vanish? State the value to which the
series converges at a jump.

\problempart[2 points] Find the Fourier coefficients of
$q(t-T/8)$ using the time-shift property. Explain what changes in the
coefficient magnitudes and phases.

\problempart[4 points] Apply $q(t)$ to the RC circuit in Problem 2, with
the periodic input present for all time. Write the Fourier series of the
periodic steady-state output using $a_k$ and $H_C(jk\omega_0)$.
Find the DC output. Explain how the circuit affects high-frequency harmonics.

\vspace*{\fill}
\begin{center}[use additional pages if needed]\end{center}

\problem{5} \textbf{Parseval, average energy, and approximation [15 points]}
% Adapted from Fall 2025 HW8 Problem 2.
Assuming other things are same, compare..

Consider the following three real, $T$-periodic signals:
\[
s(t)=A\cos(\omega_0t),\qquad b(t)=2q(t)-A,
\]
and the sawtooth defined over one period by
\[
r(t)=\frac{2A}{T}t,\qquad -T/2\le t<T/2,
\qquad r(t+T)=r(t).
\]
Here $q(t)$ is the pulse train in Problem 4, so $b(t)$ takes values $+A$
and $-A$ for equal durations.

\problempart[4 points] Find the complex Fourier coefficients of $s(t)$,
$b(t)$, and $r(t)$. Use Problem 4 for $b(t)$; for $r(t)$, show the
coefficient integral (you may adapt your sawtooth calculation from Pset 6).
Check the DC term and conjugate symmetry for each signal.

\problempart[5 points] Use the Fourier series form of Parseval's theorem
to compute the average energy over a period (power) of each signal:
\[
E_{\mathrm{avg},x}=\frac1T\int_{t_0}^{t_0+T}|x(t)|^2\,dt
   =\sum_{k=-\infty}^{\infty}|c_k|^2.
\]
Rank the signals by this quantity. You may use
\[
\sum_{k=1}^{\infty}\frac1{k^2}=\frac{\pi^2}{6},\qquad
\sum_{\substack{k\ge1\\k\ \mathrm{odd}}}\frac1{k^2}=\frac{\pi^2}{8}.
\]

\problempart[3 points] For each signal, find its energy over one period
and state whether its total energy over $(-\infty,\infty)$ is finite.
Relate the energy over one period to the average energy over a period.

\problempart[3 points] Let $d_k$ be the coefficients of $b(t)$, and define
$b_N(t)=\sum_{k=-N}^{N}d_k e^{jk\omega_0t}$. Use Parseval to compute
\[
E_N=\int_{-T/2}^{T/2}|b(t)-b_N(t)|^2\,dt
\]
for $N=1,2,3$, retaining both positive and negative harmonics. Explain any
repeated value. Does $E_N\to0$ require the Gibbs overshoot to vanish?

\vspace*{\fill}
\begin{center}[use additional pages if needed]\end{center}

\problem{6} \textbf{From Fourier series to the Fourier transform [20 points]}
Let $A>0$, $\tau>0$, and define a single, aperiodic pulse
\[
p(t)=A[u(t+\tau)-u(t-\tau)].
\]
Use $\operatorname{sinc}(z)=\sin(z)/z$ for $z\ne0$ and
$\operatorname{sinc}(0)=1$.

\problempart[4 points] Starting from the transform integral, find
$P(\omega)$. Evaluate $P(0)$ separately or by taking a limit, and relate
it to the area of the time-domain pulse.

\problempart[4 points] Sketch $P(\omega)$ and $|P(\omega)|$, labeling
the DC value and the first positive and negative zeros. If the pulse width
is doubled while its height stays fixed, how do the DC value and the first
zero locations change?

\problempart[4 points] Define $p_d(t)=p(t-t_d)$, where $t_d>0$.
Compute its transform directly by changing variables in the integral.
Compare its magnitude and phase with those of $P(\omega)$, at frequencies
where $P(\omega)\ne0$.

\problempart[4 points] Repeat $p(t)$ every $T$ seconds, with $T>2\tau$,
to form $p_T(t)$. Let $a_k(T)$ be its Fourier coefficients and
$\omega_k=2\pi k/T$. Derive the relation between $a_k(T)$ and
$P(\omega_k)$, including the case $k=0$.

\problempart[4 points] Substitute that relation into the Fourier series
synthesis equation for $p_T(t)$. With $\Delta\omega=2\pi/T$, explain
how the sum becomes the inverse Fourier transform as $T\to\infty$.
Track the factor $1/(2\pi)$ carefully. Does $a_k(T)$ itself, or the scaled
quantity $T a_k(T)$, sample the transform envelope?

\vspace*{\fill}
\begin{center}[use additional pages if needed]\end{center}

\problem{7} \textbf{Fourier transform pairs and circuit spectra [20 points]}
% Extends Fall 2025 HW8 Problem 3 with Week 8 transform-pair practice.
Use the Week 8 transform convention. Show the sifting, integration,
linearity, or inverse-transform steps that justify your answers.

\problempart[4 points] Find the transform of $\delta(t-t_0)$.
Then find the transform of the constant signal $x(t)=A$ by using the
inverse transform and sifting. Explain why the constant's forward integral
is not an ordinary convergent integral over all time.

\problempart[4 points] Find the Fourier transform of
$x(t)=e^{-at}u(t)$, with $a>0$, by evaluating the integral.
Find its value at $\omega=0$ and sketch its magnitude as a function of
$\omega$. Explain why $a>0$ is needed for this ordinary integral.

\problempart[4 points] Find the transform of
\[
x(t)=2+3\cos(4t)-2\sin(6t).
\]
Sketch the frequency impulses and label their locations and complex
weights. Use the inverse transform to check the coefficient of at least
one sinusoid.

\problempart[4 points] Find the inverse transform of
\[
X(\omega)=2\pi\delta(\omega)
 +\pi e^{j\pi/3}\delta(\omega-2)
 +\pi e^{-j\pi/3}\delta(\omega+2).
\]
Express the result as a real signal. Explain how the symmetry of these
impulse weights lets you predict that the result is real.

\problempart[4 points] Use your exponential transform pair
to find the transform of $h_C(t)$ from Problem 2.
Compare it with $H_C(j\omega)$ from Problem 3.
Sketch its magnitude, labeling the value at zero frequency and the
limit as $|\omega|\to\infty$. Does the decaying RC impulse response
contain just one frequency? Explain using its spectrum.

\vspace*{\fill}
\begin{center}[use additional pages if needed]\end{center}
\end{document}
